\documentclass[12pt,a4paper]{article}
\usepackage[utf8]{inputenc}
\usepackage[T1]{fontenc}
\usepackage[english,brazil]{babel}
\usepackage{geometry}
\geometry{top=3cm,bottom=2cm,left=3cm,right=2cm}
\usepackage{graphicx}
\usepackage{float}
\usepackage{booktabs}
\usepackage{amsmath,amssymb,amsfonts}
\usepackage{hyperref}
\hypersetup{
    colorlinks=true,
    linkcolor=blue,
    citecolor=blue,
    urlcolor=blue
}

% Define resumo environment for Portuguese abstract
\newenvironment{resumo}{%
  \otherlanguage{brazil}%
  \renewcommand{\abstractname}{Resumo}%
  \begin{abstract}%
}{%
  \end{abstract}%
}

\title{\textbf{Revisão Sistemática sobre Capacidade de Transferência Entre Idiomas de Modelos BERT}}
\author{
    Victor Hugo Santos Bitencourt \\
    Departamento de Ciências Exatas e da Terra, Campus I \\
    Universidade do Estado da Bahia (UNEB) \\
    Salvador, Bahia, Brasil. \\
    \texttt{abjurandam@gmail.com}
}
\date{}

\begin{document}

\maketitle

\begin{resumo}
Este trabalho consiste em uma revisão sistemática a sobre tal coisa, buscando investigar tal coisa. A metodologia consistiu em fazer tais coisas. Os resultados consistem em tais coisas. Conclui-se que tal coisa.
\end{resumo}

\begin{otherlanguage}{english}
\begin{abstract}
Em inglês.
\end{abstract}
\end{otherlanguage}

\section{Introdução}
No contexto do que, esse tema principal serve pra tal coisa. Qual sua importancia, relevancia, normas tecnicas? O que é esse tema principal, de acordo com tal norma tecnica? Cite alguma coisa pra nao ficar no ar. Traga o contexto do mundo, do Brasil, seja o que for.

Traga, por gentileza, a motivação do problema. Onde mora o problema? Tipo... essa área traz tais demandas, mas tal problema existe. Pode ser feito de tal jeito, mas aí tais problemas surgem. (Isso é a base pra introduzir sua solução). Traga umas definições pra já ir acostumando o leitor. Introduza termos técnicos bem de leve, os bem gerais mesmo.

Nesse cenário, essas coisas novas surgiram pra trazer tal e tal benefício. Mas tem porém, tal e tal coisa acontece. Cite algo. Assim, surge a necessidade de fazer pipipi popopo do seu jeito. ``A motivação da pesquisa surge na necessidade de investigar essas lacunas''.

No ecossistema X, sua solução se manifesta de tal forma e pode ser usada de tal forma. Isso, já vai dando sinais de sua metodologia. Contudo, tais limitações etc etc: não combina ou usa tais coisas, o que limita os ganhos. Logo, essa revisão tem como objetivo ``mapear e analisar métodos já estudados, buscando fundamentos para a construção integrada X''. Essa Revisão Sistemática de Literatura (RSL) busca identificar (liste tudo), ``práticas, mecanismos'' para esse objetivo específico. Seja completo logo.

\subsection{Motivação}
Muitos estudantes enfrentam dificuldades na formatação inicial em LaTeX. O objetivo deste modelo é fornecer uma base pronta e limpa para iniciar a escrita imediatamente.

\subsection{Estrutura do Documento}
O restante deste documento está organizado da seguinte forma:
\begin{itemize}
    \item A Seção 2 apresenta a fundamentação teórica e exemplos de inserção de imagens e fórmulas matemáticas.
    \item A Seção 3 discute a metodologia adotada e apresenta exemplos de tabelas e listas aninhadas.
    \item A Seção 4 conclui o trabalho e aponta direções futuras.
\end{itemize}

\section{Fundamentação Teórica}
A fundamentação teórica baseia-se em conceitos consolidados na literatura de engenharia de software e desenvolvimento de sistemas \cite{fowler2018refactoring}.

\subsection{Conceitos Fundamentais}
Podemos expressar relações matemáticas ou métricas de qualidade de software de forma formal:
\begin{equation}
    M = \frac{C_{\text{modificado}}}{C_{\text{total}}} \times 100
\end{equation}

Abaixo, incluímos uma imagem gerada automaticamente para ilustrar o fluxo conceitual. Graças ao uso do pacote \texttt{float} com o modificador \texttt{[H]}, a imagem permanece exatamente onde foi declarada no código-fonte, evitando deslocamentos indesejados.

\begin{figure}[H]
    \centering
    \includegraphics[width=0.8\textwidth]{figura_exemplo.png}
    \caption{Exemplo de imagem gerada por ferramenta externa e inserida no texto.}
    \label{fig:exemplo}
\end{figure}

\subsection{Padrões e Especificações}
Normas abertas são essenciais para garantir interoperabilidade na web moderna \cite{w3c2023html}.

\section{Metodologia e Análise}
A metodologia engloba as etapas de planejamento, execução e validação.

\subsection{Etapas do Processo}
Abaixo detalhamos as atividades planejadas em listas aninhadas:
\begin{enumerate}
    \item Fase de Planejamento:
          \begin{itemize}
              \item Levantamento bibliográfico inicial.
              \item Definição do escopo e limitações.
          \end{itemize}
    \item Fase de Execução:
          \begin{itemize}
              \item Desenvolvimento do protótipo ou estudo empírico.
              \item Coleta de dados e métricas.
          \end{itemize}
    \item Fase de Validação:
          \begin{itemize}
              \item Análise comparativa dos resultados.
              \item Escrita do relatório final.
          \end{itemize}
\end{enumerate}

\subsection{Comparativo de Ferramentas}
A Tabela \ref{tab:comparativo} resume as principais características das ferramentas utilizadas no projeto.

\begin{table}[H]
    \centering
    \caption{Comparativo de ferramentas e tecnologias.}
    \label{tab:comparativo}
    \begin{tabular}{lllc}
        \toprule
        \textbf{Ferramenta} & \textbf{Categoria} & \textbf{Finalidade Principal} & \textbf{Suporte} \\
        \midrule
        LaTeX / Tectonic & Editoração & Produção de documentos acadêmicos & Alto \\
        ImageMagick & Gráficos & Geração e manipulação de imagens & Alto \\
        Git & Controle de Versão & Rastreamento de alterações & Alto \\
        \bottomrule
    \end{tabular}
\end{table}

\section{Conclusão e Trabalhos Futuros}
Este template provisório fornece a estrutura necessária para iniciar a redação do TCC 1 sem preocupações excessivas com ajustes manuais de posicionamento de figuras e configuração de bibliografia. Como trabalhos futuros, sugere-se a expansão das seções teóricas e a adição de apêndices com dados brutos.

\bibliographystyle{plain}
\bibliography{references}

\end{document}
